\chapter{受体研究设计：问卷、编码与访谈}\label{ch:11}

\begin{nlkeypoints}
  \begin{enumerate}[label=(\arabic*),leftmargin=2.4em]
    \item 本章交付三套\keyw{研究设计}而非研究结果：
          面向 H7.1--H7.4 的问卷（五模块、含假设映射表）、
          面向窄通道与成本不对称命题的内容编码方案
          （含双人独立编码与一致性要求）、
          以及面向立牌活动参与者的半结构化访谈提纲。
          全部设计是供施测的工具，\keyw{本书不含任何经验数据}。
    \item 设计的三个方法论支柱：
          \keyw{学期节律对齐}（压力机制的检验窗口）、
          \keyw{不对称指标}（区分功能机制与消遣/算法解释）、
          \keyw{行为—文本—自陈三角}（防止单一来源偏差）。
    \item 伦理条款先于工具交付：
          匿名化、知情同意、不收集可识别信息、
          与第四章 E1--E3 规则衔接——
          受体研究中“研究对象是人”，
          证据分级制度之外另有全套义务。
  \end{enumerate}
\end{nlkeypoints}

第七章留下了五个待检验假设，第六章与第八章
留下了两个待检验命题，第十章留下了未测量的协商后果。
本章的任务是把这些“欠账”变成可以开工的设计。
必须在一开始就重申全书最重要的声明：
\keyw{本章不含任何测量结果}。所有量表、编码表、
提纲与示例数据都是设计（\conf{M} 级），
它们的价值在于可被批评、修订、施测与再修订
——这正是本书与“梗评”文章的最后一道分界线：
前者交付可检验的工具，后者交付不可检验的断言。
中文网络传播研究的问卷与内容分析的一般方法，
可参照彭兰的网络传播学教材与中文迷因研究的
代表性梳理\cite{penglan2007meme,caoqi2015meme}，
本章的设计与其方法传统对接，不另起炉灶。

\section{研究对象与抽样}\label{sec:11:sampling}

\subsection{总体与抽样框}

\begin{definition}[研究总体（本书界定）]\label{def:11:pop}
  研究总体为\keyw{近 12 个月内在中国大陆高校注册就读、
  且使用任一主流社交平台的在校学生}；
  目标子总体为其中\keyw{接触过奶蛙类内容}者
  （操作化：能在提示下识别奶蛙形象，
  即第六章归并标准 (C3) 的个体版本）。
\end{definition}

抽样方案的诚实版本是：\keyw{不存在可获得的完整抽样框}
（高校学生名单不可得，接触者不可枚举），
因此任何施测都是非概率抽样。设计的应对有三：
其一，\keyw{配额}：按年级（尤其区分毕业年级与非毕业年级）
与院校层次设置配额下限——毕业年级是 H7.1 与 H7.3
的关键对比组；其二，\keyw{多入口}：同好会群、
课程班级、校园公共账号三入口分开发放，
降低单入口的社群同质性；其三，\keyw{报告偏差}：
无论样本多大，结论只表述为“样本内关系”，
不得外推为“大学生总体”。
孙立平对自嘲式表达之社会心态背景的讨论
\cite{sunliping2010} 在这里进入设计动机：
若压力—自嘲的关联是广泛的社会心态，
样本内应能观察到奶蛙使用与一般自嘲表达的
\keyw{区分度}——这是第七章竞争性解释的又一条判别线。

\subsection{样本量的量级}

教学设计的参考量级：若目标是检出
“两组频率差 $\ge 15$ 个百分点”的差异
（显著性水平 0.05、功效 0.80 的双侧检验），
每组约需 130--150 份有效问卷；
考虑配额与废卷，总投放量建议不低于 600 份。
这一估算按常规两比例检验的样本量公式即可复核，
本书不在此展开推导（\conf{M} 级）。

\section{问卷设计}\label{sec:11:questionnaire}

问卷分五模块（A--E），设计原则：
每题服务于一个明确的假设或对照，
不设“你觉得奶蛙为什么火”这类开放式归因题
——归因题测到的是受访者的民间理论，不是机制本身。

\begin{table}[htbp]\centering
  \caption{问卷模块与假设映射（\conf{M} 级设计）}
  \label{tab:11:modules}
  \small
  \begin{tabularx}{\textwidth}{@{}lXX@{}}
    \toprule
    模块 & 内容 & 服务的假设 / 对照 \\ \midrule
    A 使用行为 & 各类内容的使用频率（奶蛙类 / 总娱乐类，
      按学期时段分层回填） & H7.1 的选择性上升；
      算法解释的曝光控制 \\
    B 动机与情绪 & 压力自评（PSS 风格简表）+ FoMO 量表 +
      使用情境勾选（自述 / 群聊复读 / 评论区） &
      H7.1、H7.2、H7.4、FoMO 动因 \\
    C 圈内知识 & 场景语法测验（衍生语—场景配对，
      含干扰项）+ 谱系识别 & 立场维度的客观化；
      圈内 / 圈外分组 \\
    D 自我呈现 & 主页人设自评 + 敏感期发布行为回填 &
      H7.3 的劣化内容占比与人设关联 \\
    E 元信息 & 年级、住宿类型、平台、入口来源 &
      配额核查；H7.4 的在场密度分组 \\
    \bottomrule
  \end{tabularx}
\end{table}

\subsection{示例题目}

以下给出各模块的代表题（完整量表施测前须经预试与
信度检验，\conf{M} 级）：

\begin{enumerate}[label=\textbf{Q11.\arabic*},leftmargin=4.2em,itemsep=0.3em]
  \item \textbf{（A，频率分层）}回想本学期的三个时段
        （开学初 / 期中前后 / 期末与考试周），
        你在其中刷到并\keyw{主动转发或发送}奶蛙类内容的频率
        分别是：0 次 / 每月 1--3 次 / 每周 1--2 次 /
        每周 3--5 次 / 几乎每天。
        （同题式样对“其他搞笑内容”重复一遍，作对照基线。）
  \item \textbf{（B，情境勾选）}你发送奶蛙类内容时，
        通常是为了：（可多选）表达自己的失败 / 焦虑 /
        群里有人先发了跟着发 / 纯粹觉得好笑 /
        回应别人的自嘲 / 其他\_\_。
  \item \textbf{（C，场景语法）}“无了”最适合出现的场景是：
        （单选）奶蛙出场 / 奶蛙吃瘪或惊讶 /
        奶蛙加油 / 奶蛙思考。
        （干扰项从表 \ref{tab:06:langs} 的其他场景取。）
  \item \textbf{（D，人设与发布）}考试周里，
        你是否发布过“故意做坏”的内容
        （低画质梗图、摆烂文案、自我丑化图）？
        频率与平时相比：更多 / 持平 / 更少 / 我从不发布。
\end{enumerate}

\subsection{从题目到假设的检验}

\begin{table}[htbp]\centering
  \caption{假设—题目—判别统计量对照（\conf{M} 级设计）}
  \label{tab:11:hypmap}
  \small
  \begin{tabularx}{\textwidth}{@{}lXX@{}}
    \toprule
    假设 & 所需题目 & 判别统计量 \\ \midrule
    H7.1 & Q11.1 两版频率 + B 模块压力分 &
      奶蛙类与总娱乐类的“考试周/平时”频率比之差；
      该差与压力自评的相关 \\
    H7.2 & Q11.2 情境 + （可选）行为日志 &
      情境勾选中“自述”占比的圈内 / 圈外差异 \\
    H7.3 & Q11.4 + D 模块人设自评 &
      敏感期劣化发布的组间（高/低人设）差异 \\
    H7.4 & E 模块住宿类型 + A 模块衰减回填 &
      集体 / 非集体居住组的持续使用时长比较 \\
    FoMO & B 模块 FoMO 量表 &
      FoMO 得分对入会 / 转发行为的增量解释力 \\
    \bottomrule
  \end{tabularx}
\end{table}

注意 Q11.2 的局限：情境勾选测到的仍是
\keyw{自陈动机}，受社会赞许性影响
（“跟着发”比“表达焦虑”更易勾选）。
H7.2 的强检验需要\keyw{行为数据}
（自愿提供的聊天记录片段，经脱敏），
问卷只提供弱证据——这一主次关系写进设计，
防止施测者把自陈当行为。

\section{内容编码方案}\label{sec:11:coding}

第六章命题 \ref{prop:06:channel}（窄通道）与第八章
命题 \ref{prop:08:asym}（成本不对称）需要内容编码。
方案如下（\conf{M} 级设计）：

\begin{enumerate}[label=\arabic*.]
  \item \textbf{样本框}：选定平台、选定时间窗
        （建议覆盖一次期中考试周），以固定检索式
        抓取 / 人工登记奶蛙相关内容，
        记录登记规则使他人可重复；
  \item \textbf{特征清单}：轮廓特征（黄色主体、大笑姿态、
        眼睛形态）与细节特征（配色变化、道具、
        娘化 / 物化元素、背景元素）分列，
        每条内容按清单二值编码；
  \item \textbf{窄通道检验}：比较轮廓特征保有率与
        细节特征保有率的离散系数
        （第六章预测：前者显著更小）；
  \item \textbf{成本不对称检验}：对每条内容另编
        “形式新颖度”（是否改动音轨、结尾套路、时长结构）
        与“内容新颖度”（题材、场景、跨界对象），
        比较两列的多样性指数随时间的变化
        （第八章预测：内容升、形式不升）；
  \item \textbf{质量控制}：双人独立编码，
        报告每项特征的 Krippendorff $\alpha$
        或 Cohen's $\kappa$；
        $\alpha < 0.67$ 的特征项作废重定义，
        不得带病入表；
  \item \textbf{“故意劣化”编码}（服务 H7.3）：
        只编码可观察标记（刻意低分辨率、错位排版、
        拒绝修正的手绘痕迹），
        \keyw{不}编码“看起来做得很差”——
        后者无法与真的做差区分（第七章已警告）。
\end{enumerate}

\section{访谈设计}\label{sec:11:interview}

访谈服务于问卷够不到的三类问题：
机制的\keyw{过程描述}（一次转发是怎么决定的）、
协商的\keyw{多方视角}（第十章遗留问题）、
以及圈内知识的\keyw{习得史}。
设计为半结构化，提纲如下（\conf{M} 级）：

\begin{enumerate}[label=I11.\arabic*,leftmargin=3.2em,itemsep=0.2em]
  \item 请回忆你第一次见到奶蛙的场合。当时你在做什么？
  \item 讲一次你\keyw{没有}转发但差点转发的经历——
        是什么让你停下的？
  \item 你的群聊里，谁是发奶蛙最多的人？
        你怎么看他 / 她？
  \item （立牌活动参与者）从决定投放到立牌被带走，
        哪个时刻你的想法变了？
  \item 你觉得学校管理方怎么看这些立牌？
        他们说得有道理的部分是什么？
  \item 如果明天起所有平台都屏蔽了奶蛙，
        你会做什么？
\end{enumerate}

设计要点：I11.2 用“没转发的瞬间”绕开动机的
事后美化（行为发生在 30 秒内，自陈发生在事后，
两者天然有落差）；I11.5 是第十章“管理方缺席”的
直接回应——请受访者替对方陈述理由，
可部分逼近协商的全景；I11.6 是 M5 的个体版
试探（压制下的行为意向），其结果只能作假设
生成材料，不作检验。

\section{伦理条款}\label{sec:11:ethics}

受体研究的伦理义务独立于、且优先于方法设计：

\begin{enumerate}[label=(\arabic*),leftmargin=2.4em]
  \item \textbf{知情同意}：施测前告知研究目的、
        数据用途、退出权利；未成年人不在总体内，
        若问卷中出现则终止作答；
  \item \textbf{最小化收集}：不收集姓名、学号、
        精确院系等可识别信息；年级用区间，
        院校由受访者自选披露层级；
  \item \textbf{脱敏与销毁}：行为数据（聊天记录片段）
        须去除对话方身份后提交，
        研究结束后按预先声明的期限销毁；
  \item \textbf{访谈伦理}：访谈录音须单独同意；
        涉及校方管理者的叙述时，遵守第四章 E1 规则
        ——以\keyw{角色}而非\keyw{个人}为宾语记录；
  \item \textbf{结果回馈}：向参与者社群（如同好会）
        公开去标识的结果摘要，回报其协作。
\end{enumerate}

\begin{nlmethod}[示例数据的使用纪律]
  为训练分析流程，本章配套示例数据表
  （虚构！\conf{M} 级）：40 个个案、
  六个变量（年级、住宿、考试周频率比、压力分、
  FoMO 分、圈内测验得分），文件位于
  \fname{figures/sample\_data\_example.csv}。
  纪律：(1) 该表数字为本书构造，
  \keyw{任何场合不得引用为其代表真实分布}；
  (2) 用途限于练习 \ref{tab:11:hypmap} 的统计流程；
  (3) 练习结论与真实检验无关。
  示例表的存在本身是一个教学装置：
  让读者在拿到真实数据之前，
  先见过“数据$\to$判别统计量$\to$假设裁决”的完整链路。
\end{nlmethod}

\section{本章小结}\label{sec:11:summary}

本章交付了三套设计：问卷（五模块、示例题、
假设—题目—统计量对照）、内容编码（特征清单、
窄通道与成本不对称的检验步骤、双人编码一致性要求、
“故意劣化”的标记式编码）、访谈（六题提纲、
过程视角与多方视角的设计要点），以及伦理条款五条。
设计的三个支柱——学期节律对齐、不对称指标、
行为—文本—自陈三角——分别是第七章
竞争性解释框架、第一章方法承诺与第四章考据伦理
在受体研究中的落位。至此，全书前十章欠下的
检验债务都有了对应的工具。遗留问题只有一个，
也是最大的一个：\keyw{工具尚未上过战场}。
第十二章在方法论层面收束全书，
其中“可复现性清单”即为这些工具的通用操作规程。

\begin{nlexercises}
  \item Q11.1 用“回想三个时段”回填频率，
        这种回溯式测量有系统性偏差
        （近因效应、情绪一致性回忆）。
        设计一个前瞻式的替代测量
        （提示：固定样本日记法），
        并比较两者在成本与效度上的取舍。
  \item 编码方案第 6 条拒绝编码“看起来做得很差”。
        有人反驳：“故意劣化的辨识恰恰靠整体观感，
        拆成标记就丢了。”请为双方各写一条最强理由，
        然后给出你的裁决与折中方案。
  \item I11.6（“如果平台全部屏蔽”）在设计审查中
        可能被质疑“引导受访者进入受害想象”。
        修改该题的措辞与位置（相对其他题），
        使诱导最小化；若修改后题目失去分析价值，
        说明为什么。
\end{nlexercises}

\nlkeywords 受体研究、抽样框、配额抽样、学期节律对齐、
不对称指标、自陈动机、行为数据、内容编码、
编码者一致性、半结构化访谈、知情同意、最小化收集、
示例数据
